\section{Vertex visits and finite chord means}\label{c01-c1-the-two-leg-mass}

The analytic argument has produced a partition function for polygons.
We now insert a visited vertex into a path. The conversion is a finite
exploration identity, valid before any infinite-volume limit. Its
positive expansion will also justify the deletion comparison needed
for the later transfer from arches to bridges.

A finite vertex set $D$ is \emph{filled} if every lattice vertex
strictly inside a simple polygon in its induced honeycomb graph also
belongs to $D$. Boundary ports are midpoints of edges with one endpoint
in $D$ and one outside. An excursion joins two distinct boundary ports
and visits distinct vertices of $D$. We order its endpoints, so
reversal gives a second excursion. For $v\in D$, define
\[
 T_D(v)=\sum_{\gamma\ni v}x_*^{|\gamma|}e^{i\sigma W(\gamma)},
 \qquad \sigma=3/8.
\]
Reversal conjugates each term, so $T_D(v)$ is real. Its individual
terms need not have positive real part in a general filled set.
For a face $f$ incident to $v$, let $Z_D(f)$ be the sum over systems
of vertex-disjoint polygons in $D$ enclosing its center, with weight
$2x_*^{|P|}$ per polygon and with the empty system included.

\begin{proposition}[Finite visit recursion]\label{prop:pv:visits}
For every finite filled vertex set $D$ and $v\in D$,
\begin{align}
 T_D(v)=3
 &+\sum_{\substack{P\subset D\\v\in P}}
 2x_*^{|P|-1}\big[-\cos\big(\sigma(2\pi-\tau_v(P))\big)\big]
 \label{eq:pv:visit-recursion}\\
 &+\sum_{\substack{P\subset D\\v\in P^\circ}}
 2x_*^{|P|}T_{P^\circ}(v).\notag
\end{align}
Here $P^\circ$ is the strict interior vertex set and
$\tau_v(P)\in\{\pi/3,-\pi/3\}$ is the counterclockwise turn at $v$.
Every coefficient on the right is positive. Thus $T_D(v)\ge3$,
it is monotone under inclusion of filled sets containing $v$, and
there are absolute constants $c,C>0$ such that
\[
 c\sum_{f\sim v}Z_D(f)\le T_D(v)\le C\sum_{f\sim v}Z_D(f).
\]
If $D$ is cut out by admissible tiling lines in a convex polygon and
$A_D(v)$ is the positive ordered excursion weight through $v$, then
\[
 \cos(3\pi/8)A_D(v)\le T_D(v)\le A_D(v).
\]
\end{proposition}

\begin{proof}
At a fresh vertex reached along one edge, take either nonbacktracking
continuation with weight $x_*e^{i\sigma\pi/3}$ or
$x_*e^{-i\sigma\pi/3}$. Their sum is one. Stop at an exit port or just
before choosing a continuation at the first repeated vertex. The
exploration tree is finite and replacing a node by its children
preserves total weight. A prefix stopped on first arrival at a vertex
has not yet taken its weight or its turn.

Start at a boundary port. A first repeat consists of a simple stem
followed by a polygon attached at its end. The stem is outside the
polygon: otherwise the outside neighbor of its starting port would
be inside that polygon, contradicting filling. The two orientations
of the polygon add turns $\pm4\pi/3$ relative to the stem. Their
phase sum is $2\cos(4\pi\sigma/3)=0$. This pairing also cancels after
restricting to histories that have visited $v$, because reversal
preserves their vertices. Continuing a first-arrival prefix at $v$
has total terminal weight one. Hence the sum of first-arrival prefixes
equals the excursion weight through $v$. Summing starting ports gives
$T_D(v)$.

Now start at $v$ along each of its three edges, with initial weight one
and with neither weight nor turn at $v$. Stop at an exit, a return to
$v$, or the first repeat of another vertex. The three explorations
have total weight $3$. Their exit terms reverse the prefixes just
counted, so their sum is the conjugate of $T_D(v)$ and hence $T_D(v)$.
A return to $v$ traces a polygon $P$ through it. Its turn omits exactly
$\tau_v(P)$ from the total $2\pi$, and its weight omits exactly one
factor $x_*$. The two orientations therefore contribute
\[
 2x_*^{|P|-1}\cos\big(\sigma(2\pi-\tau_v(P))\big).
\]

For a first repeat away from $v$, let $P$ be the attached polygon. If
$v$ is outside $P$, exterior-stem cancellation applies. If $v$ is
inside, the unused edge at the attachment points into the polygon.
Arriving from the stem replaces the counterclockwise turn $-\pi/3$
there by $+\pi/3$. The two traversals add turns $\pm8\pi/3$,
both with phase $-1$, and contribute $-2x_*^{|P|}$ times the stem
weight. For fixed $P$, these stems are exactly paths from $v$, with
no weight or turn there, to boundary ports of $P^\circ$. Filling gives
$P^\circ\subset D$; planarity shows that an edge leaving its interior
ends on $P$. Completing the last half-edge therefore introduces no
extra weighted vertex. Prefix reversal in this smaller filled set
gives stem sum $T_{P^\circ}(v)$. The terminal balance is
\[
 3=T_D(v)+\sum_{P\ni v}2x_*^{|P|-1}
 \cos\big(\sigma(2\pi-\tau_v(P))\big)
 -\sum_{v\in P^\circ}2x_*^{|P|}T_{P^\circ}(v),
\]
which proves the recursion.

If $P$ encloses one of the three faces incident to $v$, its turn there
is $+\pi/3$; if it encloses two, the turn is $-\pi/3$. The omitted-turn
angles are $5\pi/3$ and $7\pi/3$. Set
\[
 c_1=-\cos(5\pi/8)>0,\qquad c_2=-\cos(7\pi/8)>0.
\]
Iterate the recursion. It terminates since a strict interior has fewer
vertices than its container. The result sums outer-to-inner chains
of disjoint polygons strictly enclosing $v$, with factor $2x_*^{|P|}$
for each. In the final interior $D'$ the remaining multiplier is
\[
 3+\frac2{x_*}\sum_{\substack{Q\subset D'\\v\in Q}}
 c_{m_Q}x_*^{|Q|},\qquad m_Q\in\{1,2\}.
\]
This positive expansion proves both positivity and filled-set
monotonicity: inclusion only adds eligible chains and through-polygons.

The same strict chains enumerate $\sum_{f\sim v}Z_D(f)$. A polygon
strictly enclosing $v$ encloses all three incident face centers. A
polygon through $v$ encloses exactly $m_Q$ of them. Conversely, a
polygon enclosing an incident face has $v$ in its interior or on its
boundary, because the open face does not meet the polygon and
approaches $v$. A vertex-disjoint system has at most one polygon
through $v$, and it must be innermost. Thus its final multiplier is
\[
 3+2\sum_{\substack{Q\subset D'\\v\in Q}}m_Qx_*^{|Q|}.
\]
The two ratios $c_m/(x_*m)$ are fixed positive numbers. Termwise
comparison proves the incident-face bounds.

Finally close a convex-domain excursion by the complementary
counterclockwise boundary arc. Its turn is the lifted change of
outward boundary normal from the initial port to the terminal one,
minus $\pi$. That change lies in $[0,2\pi]$, so $|W|\le\pi$.
Each phased real part is therefore between $\cos(3\pi/8)$ and one
times its positive weight. Summation proves the last assertion.
\end{proof}

The proposition estimates visits without choosing the wall on which
the excursion ends. We first apply it to hexagon chords and mixed
strip exits; only afterward do we isolate bridges.

